% GENERATED FILE. Edit canonical lessons, then run npm run generate:books.
% canonical-source-hash: fnv1a64:95986e3a7864c7d3
% canonical-lessons: TE-C241-mingu, TE-C241-koruku, TE-C241-pilcu, TE-C241-tummu, TE-C241-vanti-cesuko

\chapter{To swallow, To bite, To breathe in, To sneeze, To vomit}
\label{ch:te-to-swallow-to-bite-to-breathe-in-to-sneeze-to-vomit}
\hlchaptermodality{241}

\begin{chapteropening}
\textbf{By the end of this chapter:} \emph{I can say to swallow, to bite, to breathe in, to sneeze and to vomit.}

Complete the last lesson of chapter 241: I can say to swallow, to bite, to breathe in, to sneeze and to vomit.
\end{chapteropening}

\section[\texorpdfstring{\te{మింగు} (miṅgu)}{మింగు (miṅgu)}]{\te{మింగు} (miṅgu) — to swallow}

\label{lesson:TE-C241-mingu}



\begin{quote}
Before the new one: say the Telugu for to push, then the Telugu for to throw.
\end{quote}

\subsection*{You'll want to know: \te{మింగు}}
\textbf{\te{మింగు}} — \emph{miṅgu} — \textquotedblleft{}to swallow\textquotedblright{}.

\begin{grammarlens}[title={What you've built}]
One more word for this chapter.
\end{grammarlens}

\subsection*{Your turn}
\begin{itemize}
\raggedright
  \item \emph{Say it:} \emph{miṅgu}
  \item \emph{Say it:} \emph{miṅgu}, once more
  \item \emph{Say it:} say \emph{miṅgu}
  \item \emph{Recall:} say the Telugu for to push, then the Telugu for to throw, then say \emph{miṅgu} again
\end{itemize}

\begin{tcolorbox}[breakable,colback=teal!4,colframe=teal!35!black,title={Before you move on}]
What does \te{మింగు} mean? (\textquotedblleft{}To swallow\textquotedblright{}.) Say it once more.
\end{tcolorbox}
\section[\texorpdfstring{\te{కొరుకు} (koruku)}{కొరుకు (koruku)}]{\te{కొరుకు} (koruku) — to bite}

\label{lesson:TE-C241-koruku}



\begin{quote}
Before the new one: say the Telugu for to throw, then the Telugu for to swallow.
\end{quote}

\subsection*{You'll want to know: \te{కొరుకు}}
\textbf{\te{కొరుకు}} — \emph{koruku} — \textquotedblleft{}to bite\textquotedblright{}.

\begin{grammarlens}[title={What you've built}]
One more word for this chapter.
\end{grammarlens}

\subsection*{Your turn}
\begin{itemize}
\raggedright
  \item \emph{Say it:} \emph{koruku}
  \item \emph{Say it:} \emph{koruku}, once more
  \item \emph{Say it:} say \emph{koruku}
  \item \emph{Recall:} say the Telugu for to throw, then the Telugu for to swallow, then say \emph{koruku} again
\end{itemize}

\begin{tcolorbox}[breakable,colback=teal!4,colframe=teal!35!black,title={Before you move on}]
What does \te{కొరుకు} mean? (\textquotedblleft{}To bite\textquotedblright{}.) Say it once more.
\end{tcolorbox}
\section[\texorpdfstring{\te{పీల్చు} (pīlcu)}{పీల్చు (pīlcu)}]{\te{పీల్చు} (pīlcu) — to breathe in, to inhale}

\label{lesson:TE-C241-pilcu}



\begin{quote}
Before the new one: say the Telugu for to swallow, then the Telugu for to bite.
\end{quote}

\subsection*{You'll want to know: \te{పీల్చు}}
\textbf{\te{పీల్చు}} — \emph{pīlcu} — \textquotedblleft{}to breathe in, to inhale\textquotedblright{}.

\begin{grammarlens}[title={What you've built}]
One more word for this chapter.
\end{grammarlens}

\subsection*{Your turn}
\begin{itemize}
\raggedright
  \item \emph{Say it:} \emph{pīlcu}
  \item \emph{Say it:} \emph{pīlcu}, once more
  \item \emph{Say it:} say \emph{pīlcu}
  \item \emph{Recall:} say the Telugu for to swallow, then the Telugu for to bite, then say \emph{pīlcu} again
\end{itemize}

\begin{tcolorbox}[breakable,colback=teal!4,colframe=teal!35!black,title={Before you move on}]
What does \te{పీల్చు} mean? (\textquotedblleft{}To breathe in, to inhale\textquotedblright{}.) Say it once more.
\end{tcolorbox}
\section[\texorpdfstring{\te{తుమ్ము} (tummu)}{తుమ్ము (tummu)}]{\te{తుమ్ము} (tummu) — to sneeze}

\label{lesson:TE-C241-tummu}



\begin{quote}
Before the new one: say the Telugu for to bite, then the Telugu for to breathe in.
\end{quote}

\subsection*{You'll want to know: \te{తుమ్ము}}
\textbf{\te{తుమ్ము}} — \emph{tummu} — \textquotedblleft{}to sneeze\textquotedblright{}.

\begin{grammarlens}[title={What you've built}]
One more word for this chapter.
\end{grammarlens}

\subsection*{Your turn}
\begin{itemize}
\raggedright
  \item \emph{Say it:} \emph{tummu}
  \item \emph{Say it:} \emph{tummu}, once more
  \item \emph{Say it:} say \emph{tummu}
  \item \emph{Recall:} say the Telugu for to bite, then the Telugu for to breathe in, then say \emph{tummu} again
\end{itemize}

\begin{tcolorbox}[breakable,colback=teal!4,colframe=teal!35!black,title={Before you move on}]
What does \te{తుమ్ము} mean? (\textquotedblleft{}To sneeze\textquotedblright{}.) Say it once more.
\end{tcolorbox}
\section[\texorpdfstring{\te{వాంతి} \te{చేసుకో} (vānti cēsukō)}{వాంతి చేసుకో (vānti cēsukō)}]{\te{వాంతి} \te{చేసుకో} (vānti cēsukō) — to vomit}

\label{lesson:TE-C241-vanti-cesuko}



\begin{quote}
Before the new one: say the Telugu for to breathe in, then the Telugu for to sneeze.
\end{quote}

\subsection*{You'll want to know: \te{వాంతి} \te{చేసుకో}}
\textbf{\te{వాంతి} \te{చేసుకో}} — \emph{vānti cēsukō} — \textquotedblleft{}to vomit\textquotedblright{}.

\begin{grammarlens}[title={What you've built}]
One more word for this chapter.
\end{grammarlens}

\subsection*{Your turn}
\begin{itemize}
\raggedright
  \item \emph{Say it:} \emph{vānti cēsukō}
  \item \emph{Say it:} \emph{vānti cēsukō}, once more
  \item \emph{Say it:} say \emph{vānti cēsukō}
  \item \emph{Recall:} say the Telugu for to breathe in, then the Telugu for to sneeze, then say \emph{vānti cēsukō} again
\end{itemize}

\begin{tcolorbox}[breakable,colback=teal!4,colframe=teal!35!black,title={Before you move on}]
What does \te{వాంతి} \te{చేసుకో} mean? (\textquotedblleft{}To vomit\textquotedblright{}.) Say it once more.
\end{tcolorbox}
